\chapter{Introduction}

\section{Background and Motivation}

In contemporary society, maintaining optimal health and fitness has become increasingly challenging due to busy lifestyles, dietary complexities, and the abundance of conflicting health information. The World Health Organization reports that lifestyle-related diseases are among the leading causes of mortality worldwide, emphasizing the critical need for accessible health management tools. While numerous health and fitness applications exist in the market, many lack comprehensive features that adequately address personalized health needs, dietary restrictions, and medical conditions.

The motivation for developing NutriFit stems from the identified gap in existing health management solutions. Current applications often provide generic recommendations without considering individual health profiles, medical conditions, or food allergens. Furthermore, many solutions fail to integrate nutrition tracking, exercise management, and progress monitoring in a cohesive, user-friendly platform. This fragmentation leads to poor user engagement and inconsistent health tracking.

The increasing prevalence of lifestyle-related health conditions, coupled with growing health consciousness among individuals, underscores the importance of accessible tools that facilitate informed health decisions. NutriFit addresses this need by providing evidence-based recommendations while ensuring user safety through intelligent validation and personalized constraints based on established nutritional science principles such as the Harris-Benedict equation for calorie calculation.

\section{Problem Statement}

The primary problem addressed by this thesis is the lack of comprehensive, personalized, and safety-aware health tracking applications that integrate multiple aspects of health management. Specifically, the following challenges exist in current solutions:

\begin{enumerate}
    \item Absence of personalized diet and workout plan generation based on individual health profiles including age, weight, height, gender, activity level, and health goals
    \item Inadequate safety mechanisms for identifying food allergens and exercise restrictions based on medical conditions
    \item Lack of integrated progress tracking features that motivate consistent user engagement
    \item Poor user experience across different devices due to non-responsive design
    \item Insufficient validation of user inputs leading to inaccurate health data
\end{enumerate}

These limitations result in users either abandoning health tracking altogether or using multiple disconnected applications, leading to fragmented data and reduced effectiveness in achieving health goals.

\section{Objectives of the Study}

The primary objectives of this thesis are:

\begin{enumerate}
    \item To design and implement a comprehensive web-based health and fitness tracking application that addresses the limitations of existing solutions
    \item To develop intelligent algorithms for generating personalized diet and workout plans based on individual health profiles
    \item To implement robust safety mechanisms for identifying food allergens and exercise restrictions
    \item To create an intuitive, responsive user interface that encourages consistent health tracking across multiple device types
    \item To validate the system through comprehensive testing ensuring reliability, performance, and usability
\end{enumerate}

\section{Scope of the Study}

This thesis encompasses the complete development lifecycle of the NutriFit application, including:

\textbf{Included in Scope:}
\begin{itemize}
    \item Requirements analysis and system design
    \item Database schema design with 11 interconnected tables
    \item Backend API development using Next.js API routes
    \item Frontend development with React and TypeScript
    \item User authentication and authorization using JWT tokens
    \item Implementation of core features: health profiles, meal logging, exercise tracking, goal setting
    \item Advanced features: favorites system, water intake tracking, recipe suggestions, reminders
    \item Comprehensive testing at unit, integration, and system levels
    \item Responsive design for mobile, tablet, and desktop devices
\end{itemize}

\textbf{Excluded from Scope:}
\begin{itemize}
    \item Integration with external fitness tracking devices or health apps
    \item Native mobile applications for iOS and Android
    \item Real-time collaboration features or social networking capabilities
    \item Payment processing or subscription management
    \item Advanced machine learning models for predictive analytics
\end{itemize}

\section{Methodology Overview}

The development of NutriFit followed a systematic approach combining agile methodologies with structured software engineering practices:

\begin{enumerate}
    \item \textbf{Requirements Gathering}: Identified functional and non-functional requirements through analysis of existing health tracking applications and user needs
    \item \textbf{System Design}: Designed a three-tier architecture with clear separation of concerns between presentation, application, and data layers
    \item \textbf{Database Design}: Created a normalized database schema using Prisma ORM with SQLite, ensuring data integrity through proper relationships and constraints
    \item \textbf{Iterative Development}: Implemented features incrementally, starting with core functionality (authentication, health profiles) and progressively adding advanced features
    \item \textbf{Testing}: Conducted unit tests for core calculations, integration tests for API endpoints, and system tests for complete user workflows
    \item \textbf{Validation}: Evaluated performance metrics including page load times, API response times, and browser compatibility
\end{enumerate}

The technology stack was selected based on modern web development best practices, prioritizing type safety (TypeScript), developer productivity (Next.js), and maintainability (Prisma ORM).

\section{Structure of the Thesis}

This thesis is organized into four chapters following the structure recommended by the supervisor:

\textbf{Chapter 1: Introduction} outlines the background, motivation, problem statement, objectives, scope, and methodology of the project.

\textbf{Chapter 2: User Documentation} describes how the system is used, including installation procedures, system requirements, user interface overview, and detailed functional descriptions of all features.

\textbf{Chapter 3: Developer Documentation} explains the internal design, system architecture, database schema, technologies used, application components, security mechanisms, and testing methodology.

\textbf{Chapter 4: Summary} concludes the work by summarizing development overview, system functionality, testing results, achievements, and possible future improvements.

Additionally, the thesis includes appendices containing key code samples and complete database schema, followed by bibliography and lists of figures, tables, codes, and abbreviations.
